%%%%%%%%%%%%%%%%%%%%%%%%%%%%%%%%%%%%%%%%%%%%%%%%%%%%%%%%%%%%%%%%%%%%%%%
% Template Skripsi - Fakultas Ilmu Komputer Universitas Brawijaya
% Berbasis Harvard Anglia referencing style
%
% Cara penggunaan:
%   1. Isi data diri pada bagian "Fill your data here"
%   2. Tulis konten bab pada file bab1.tex, bab2.tex, dst.
%   3. Tambahkan referensi pada daftar-pustaka.bib
%   4. Build dengan XeLaTeX (latexmk -xelatex)
%%%%%%%%%%%%%%%%%%%%%%%%%%%%%%%%%%%%%%%%%%%%%%%%%%%%%%%%%%%%%%%%%%%%%%%
\documentclass{skripsi}

% ------------------------------------------------------------
% Add your own needed package here
% ------------------------------------------------------------
\usepackage[titles]{tocloft}
\renewcommand\cftfigpresnum{Gambar\  }
\renewcommand\cfttabpresnum{Tabel\   }

\usepackage{hyperref}
\newlength{\mylenf}
\settowidth{\mylenf}{\cftfigpresnum}
\setlength{\cftfignumwidth}{\dimexpr\mylenf+2em}
\setlength{\cfttabnumwidth}{\dimexpr\mylenf+2em}

\usepackage[center,labelfont=bf,font=bf,labelsep=space]{caption}
\usepackage{subcaption}
\usepackage{float}
\usepackage{csquotes}
\usepackage{enumitem}
\usepackage{titlesec}

\titleformat{\paragraph}
{\normalfont\normalsize\bfseries}{\theparagraph}{1em}{}
\titlespacing*{\paragraph}
{0pt}{3.25ex plus 1ex minus .2ex}{1.5ex plus .2ex}

\usepackage[all]{nowidow}
\usepackage{ragged2e}
\usepackage{longtable, makecell}
\newcolumntype{L}[1]{>{\RaggedRight}p{#1}}
\renewcommand\theadfont{\small\bfseries}
\setcellgapes{3pt}
\setlength{\parskip}{0.25em}
\emergencystretch=1em
\usepackage{array}
\newcolumntype{P}[1]{>{\raggedright\let\newline\\\arraybackslash\hspace{0pt}}p{#1}}

% -----------------------------------------------------------------
% Custom source code style
% -----------------------------------------------------------------
\usepackage{listings}
\usepackage[newfloat,chapter]{minted}

\newenvironment{code}{\captionsetup{type=listing}}{}
\SetupFloatingEnvironment{listing}{name=Tabel Kode}

\usepackage{tcolorbox}
\tcbuselibrary{listings,minted,skins,breakable}
\usetikzlibrary{decorations.pathmorphing}
\usetikzlibrary{decorations.pathreplacing}
\newtcblisting{ignasicblock}[1][]{%
  breakable,
  colback=white,
  colframe=black,
  colbacktitle=white,
  sharp corners,
  enhanced,
  listing engine=minted,
  listing only,
  left=10mm,
  title=Source Code,
  halign title=center,
  overlay={\draw[line width=.5mm] ([xshift=8mm]frame.south west)
    -- ([xshift=8mm]frame.north west);
    \node[right] at (title.west) {No};},
  minted style=colorful,
  minted language=Python,
  minted options={%
    linenos=true,
    fontsize=\footnotesize,
    numbersep=6mm,
    texcl=true,
    breaklines=true,
    autogobble=true},
  coltitle=black,
  #1
}

\newcommand{\drawbrace}[5][]{%
  \draw [decorate,decoration={brace,amplitude=5pt},blue!75!black,line width=1pt]
    ([xshift=#2,yshift=-3.3mm-#3*12pt]interior.north west)
    -- ([xshift=#2,yshift=-2.7mm-#4*12pt]interior.north west)
    node [align=center,right=10pt,midway,circle,draw=blue!50,fill=blue!5,text=black,
      font=\sffamily\small,#1] {#5};
}

\newcommand{\drawline}[5][]{%
  \draw [blue!75!black,line width=1pt]
    ([xshift=#2,yshift=-3mm+6pt-#4*12pt]interior.north west)
    -- ([xshift=#3,yshift=-3mm+6pt-#4*12pt]interior.north west)
    node [align=center,right=5pt,circle,draw=blue!50,fill=blue!5,text=black,
      font=\sffamily\small,#1] {#5};
}

\definecolor{codebg}{rgb}{0.95,0.95,0.95}
\renewcommand\theFancyVerbLine{\footnotesize\arabic{FancyVerbLine}}

% -----------------------------------------------------------------
% Font style
% Ganti path sesuai lokasi font di komputer Anda
% -----------------------------------------------------------------
\usepackage{fontspec}

\setmainfont[Path=/Users/GANTI-DENGAN-USERNAME/Library/Fonts/,
    BoldItalicFont=Calibriz.ttf,
    BoldFont      =Calibrib.ttf,
    ItalicFont    =Calibrii.ttf]{Calibri.ttf}

\setmonofont[Path=/System/Library/Fonts/Supplemental/,
    BoldItalicFont=Courier New Bold Italic.ttf,
    BoldFont      =Courier New Bold.ttf,
    ItalicFont    =Courier New Italic.ttf]{Courier New.ttf}

% -----------------------------------------------------------------
% Bibliography style
% -----------------------------------------------------------------
%%%%%%%%%%%%%%%%%%%%%%%%%%%%%%%%%%%%%%%%%%%%%%%%%%%%%%%%%%%%%%%%%%%%%%%
% This file redefine reference list to match Harvard-Angila flavor
% Great thanks for Moritz (moewe) Wemheuer
%%%%%%%%%%%%%%%%%%%%%%%%%%%%%%%%%%%%%%%%%%%%%%%%%%%%%%%%%%%%%%%%%%%%%%%

\usepackage[
  backend=biber,
  style=ext-authoryear-comp, giveninits, uniquename=init, maxbibnames=999,
  maxcitenames=2, mincitenames=1,
  urldate=long,
  innamebeforetitle, articlein=false,
]{biblatex}

\DeclareLanguageMapping{bahasa}{english}

% 1.2
\DeclareDelimFormat*{nameyeardelim}{\addcomma\space}
\DeclareDelimFormat[textcite]{nameyeardelim}{\addspace}

% 2.5
\DeclareDelimFormat{andothersdelim}{\addcomma\space}

% 2.7
\renewcommand*{\compcitedelim}{\multicitedelim}

% 2.8
\makeatletter
\AtEveryCitekey{\global\undef\cbx@lastyear}
\makeatother

% ignore 2.10 for the time being

% 4.1
\DeclareNameAlias{sortname}{family-given}
\DeclareFieldFormat{biblabeldate}{#1}
\DeclareFieldAlias{biblistlabeldate}{biblabeldate}

% 4.3
\DeclareDelimFormat{editortypedelim}{\addspace}
\DeclareDelimFormat{translatortypedelim}{\addspace}

% 4.4
\DeclareFieldFormat
  [article,inbook,incollection,inproceedings,patent,thesis,unpublished]
  {title}{#1\isdot}

% as for the year of the chapterv vs year of the book in 4.4
% I have never seen such a distinction and in general I imagine
% it would be quite hard to find out the year of a chapter
% in a collection work since only the year of the collection is given
% I imagine that in practice you would almost always end up douplicating
% the year.

% 4.8 & 5.11
\DefineBibliographyStrings{english}{
  urlfrom   = {Tersedia pada},
  urlseen   = {Diakses},
  bathesis  = {BA},
  mathesis  = {MA},
  phdthesis = {PhD},
  andothers  = {et al\adddot}
}
% Indonesian month names for urldate
\DefineBibliographyExtras{english}{%
  \protected\def\mkbibmonth#1{\ifcase#1\or
    Januari\or Februari\or Maret\or April\or Mei\or Juni\or
    Juli\or Agustus\or September\or Oktober\or November\or Desember\fi}%
}
\AtBeginDocument{%
  \DeclareFieldFormat{url}{Tersedia pada\addcolon\space\url{#1}}%
  \DeclareFieldFormat{urldate}{\mkbibbrackets{Diakses\space#1}}%
}

% 4.9
\renewcommand*{\jourvoldelim}{\addcomma\space}
\renewcommand*{\volnumdelim}{}
\DeclareFieldFormat[article,periodical]{number}{\mkbibparens{#1}}

\renewcommand*{\volnumdatedelim}{\addcomma\space}
\DeclareFieldFormat{issuedate}{#1}

% 4.12 :-(
% note that I used the new preferred resolver https://doi.org/
% instead of http://dx.doi.org/
\DeclareFieldFormat{doi}{\url{https://doi.org/#1}}

\addbibresource{daftar-pustaka.bib}

\renewcommand*{\finalnamedelim}{%
  \ifnumgreater{\value{liststop}}{2}{\finalandcomma}{}%
  \addspace\&\space}

\newcommand{\cmmnt}[1]{\ignorespaces}

% -----------------------------------------------------------------
% Fill your data here (ISI DATA ANDA DI SINI)
% -----------------------------------------------------------------
\titleskripsi{Judul Skripsi Anda Di Sini}
\fullname{Nama Lengkap Anda}
\idnum{NIM Anda}
\approvaldate{Tanggal Persetujuan}
\degree{Sarjana Komputer}
\yearsubmit{2025}
\dept{TEKNIK INFORMATIKA}
\faculty{FAKULTAS ILMU KOMPUTER}
\university{universitas brawijaya}
\city{malang}
% Data dosen pembimbing
\firstsupervisor{Nama Dosen Pembimbing I, S.T., M.T., Ph.D}
\firstsupervisornip{NIP: XXXXXXXXXXXXXXXX}
\secondsupervisor{Nama Dosen Pembimbing II, S.T., M.T., Ph.D}
\secondsupervisornip{NIP: XXXXXXXXXXXXXXXX}

% -----------------------------------------------------------------
% And so it begins
% -----------------------------------------------------------------
\begin{document}

% -----------------------------------------------------------------
% Cover
% -----------------------------------------------------------------
\cover

% -----------------------------------------------------------------
% Approval Page
% -----------------------------------------------------------------
%\preapprovalpage
\approvalpage

%-----------------------------------------------------------------
% PERNYATAAN ORISINALITAS
%-----------------------------------------------------------------
{\orisinalitas

  Saya menyatakan dengan sebenar-benarnya bahwa sepanjang pengetahuan
  saya, di dalam naskah skripsi ini tidak terdapat karya ilmiah yang
  pernah diajukan oleh orang lain untuk memperoleh gelar akademik di
  suatu perguruan tinggi, dan tidak terdapat karya atau pendapat yang
  pernah ditulis atau diterbitkan oleh orang lain, kecuali yang secara
  tertulis disitasi dalam naskah ini dan disebutkan dalam daftar
  referensi.

  Apabila ternyata didalam naskah skripsi ini dapat dibuktikan
  terdapat unsur-unsur plagiasi, saya bersedia skripsi ini digugurkan
  dan gelar akademik yang telah saya peroleh (sarjana) dibatalkan,
  serta diproses sesuai dengan peraturan perundang-undangan yang
  berlaku (UU No. 20 Tahun 2003, Pasal 25 ayat 2 dan Pasal 70).
  \vspace{1.5cm}

  \noindent
  \hspace*{8cm}Malang, Tanggal Bulan Tahun   \vspace{2.5cm}   \\

  \hspace*{7cm}\underline{Nama Lengkap Anda} \\
  \hspace*{8cm}NIM : NIM Anda
}

%-----------------------------------------------------------------
% PRAKATA
%-----------------------------------------------------------------
{\preface

  Puji syukur kehadirat Allah SWT yang telah melimpahkan rahmat,
  taufik dan hidayah-Nya sehingga skripsi yang berjudul ``\textit{Judul
  Skripsi Anda}'' ini dapat terselesaikan. Dengan selesainya penulisan
  skripsi ini, penulis ingin menyampaikan terima kasih kepada:

  \begin{singlespace}
    \begin{enumerate}
    \item{Allah SWT yang telah memberikan nikmat kesehatan, kekuatan dan
        kemudahan sehingga penulis dapat menyelesaikan skripsi ini.}
    \item{Orang tua dan seluruh keluarga besar penulis yang senantiasa
        mendoakan dan mendukung demi terselesaikannya skripsi ini.}
    \item{Bapak/Ibu Nama Dosen Pembimbing I, selaku Dosen Pembimbing I
        yang telah meluangkan waktu untuk membimbing dengan semangat dan
        sabar sehingga penulis dapat menyelesaikan skripsi ini.}
    \item{Bapak/Ibu Nama Dosen Pembimbing II, selaku Dosen Pembimbing II
        yang telah mengarahkan dan membimbing penulis dengan semangat dan
        sabar sehingga penulis dapat menyelesaikan skripsi ini.}
    \item{Seluruh civitas akademika Teknik Informatika Universitas
        Brawijaya yang telah banyak memberi bantuan dan dukungan selama
        penyelesaian skripsi ini.}
    \end{enumerate}
  \end{singlespace}

  Penulis menyadari bahwa dalam penelitian ini masih banyak
  kekurangan. Oleh karena itu, saran dan kritik yang membangun sangat
  penulis harapkan. Akhir kata penulis berharap penelitian ini dapat
  membawa manfaat bagi semua pihak yang menggunakannya.

  \vspace{0.8cm}

  \noindent
  \hspace*{8cm}Malang, Tanggal Bulan Tahun  \vspace{1.5cm} \\

  \hspace*{6.8cm}Penulis \\
  \hspace*{8cm}email@student.ub.ac.id
}

% -----------------------------------------------------------------
% ABSTRAK INDONESIA
% -----------------------------------------------------------------
{\abstractind

  \noindent
  Tulis abstrak bahasa Indonesia di sini. Abstrak berisi ringkasan
  singkat dari keseluruhan penelitian, mencakup latar belakang, tujuan,
  metode, hasil, dan kesimpulan. Panjang abstrak biasanya 150--250 kata.\\

  \noindent
  \textbf{Kata Kunci}: \emph{kata kunci pertama},
  \emph{kata kunci kedua}, \emph{kata kunci ketiga}

}

% -----------------------------------------------------------------
% ABSTRAK ENGLISH
% -----------------------------------------------------------------
{\abstracteng

  \noindent
  Write your English abstract here. The abstract should contain a brief
  summary of the entire research, including the background, objectives,
  methods, results, and conclusions. Abstract length is typically
  150--250 words.\\

  \noindent
  \textbf{Keywords}: \emph{first keyword},
  \emph{second keyword}, \emph{third keyword}

}

% -----------------------------------------------------------------
% TOCS
% -----------------------------------------------------------------
\tableofcontents
\addcontentsline{toc}{chapter}{DAFTAR ISI}
\listoftables
\addcontentsline{toc}{chapter}{DAFTAR TABEL}
\listoffigures
\addcontentsline{toc}{chapter}{DAFTAR GAMBAR}
\addcontentsline{toc}{chapter}{DAFTAR LAMPIRAN}

% -----------------------------------------------------------------
% Chapter
% -----------------------------------------------------------------
\selectlanguage{bahasa}\clearpage\pagenumbering{arabic}\setcounter{page}{1}

%%%%%%%%%%%%%%%%%%%%%%%%%%%%%%%%%%%%%%%%%%%%%%%%%%%%%%%%%%%%%%%%%%%%%%%
% BAB 1 - PENDAHULUAN
%%%%%%%%%%%%%%%%%%%%%%%%%%%%%%%%%%%%%%%%%%%%%%%%%%%%%%%%%%%%%%%%%%%%%%%

\mychapter{1}{BAB 1 PENDAHULUAN}

\section{Latar Belakang}

% Tulis latar belakang penelitian di sini.
% Contoh sitasi: \parencite{contoh-artikel}
% Contoh sitasi dalam kalimat: \textcite{contoh-buku}

Lorem ipsum dolor sit amet, consectetur adipiscing elit.

\section{Rumusan Masalah}

% Tulis rumusan masalah di sini.

Berdasarkan latar belakang di atas, rumusan masalah penelitian ini adalah:
\begin{enumerate}
  \item Pertanyaan penelitian pertama?
  \item Pertanyaan penelitian kedua?
\end{enumerate}

\section{Tujuan Penelitian}

% Tulis tujuan penelitian di sini.

Tujuan dari penelitian ini adalah:
\begin{enumerate}
  \item Tujuan pertama.
  \item Tujuan kedua.
\end{enumerate}

\section{Manfaat Penelitian}

% Tulis manfaat penelitian di sini.

\section{Batasan Masalah}

% Tulis batasan masalah di sini.

\section{Sistematika Pembahasan}

\noindent
\textbf{BAB I : PENDAHULUAN}

Bab ini berisi latar belakang, rumusan masalah, tujuan, manfaat,
batasan masalah, dan sistematika pembahasan.

\noindent
\textbf{BAB II : LANDASAN TEORI}

Bab ini berisi tinjauan pustaka dan landasan teori yang mendukung
penelitian.

\noindent
\textbf{BAB III : METODOLOGI}

Bab ini menjelaskan metodologi yang digunakan dalam penelitian.

\noindent
\textbf{BAB IV : HASIL DAN PEMBAHASAN}

Bab ini memaparkan hasil penelitian beserta pembahasannya.

\noindent
\textbf{BAB V : KESIMPULAN DAN SARAN}

Bab ini berisi kesimpulan dan saran berdasarkan hasil penelitian.

%%%%%%%%%%%%%%%%%%%%%%%%%%%%%%%%%%%%%%%%%%%%%%%%%%%%%%%%%%%%%%%%%%%%%%%
% BAB 2 - LANDASAN TEORI
%%%%%%%%%%%%%%%%%%%%%%%%%%%%%%%%%%%%%%%%%%%%%%%%%%%%%%%%%%%%%%%%%%%%%%%

\mychapter{2}{BAB 2 LANDASAN TEORI}

\section{Tinjauan Pustaka}

% Tulis tinjauan pustaka di sini.
% Gunakan \textcite{key} untuk sitasi dalam kalimat.
% Gunakan \parencite{key} untuk sitasi dalam kurung.

\section{Landasan Teori}

\subsection{Subbab Pertama}
\label{subsec:subbab-pertama}

% Tulis konten di sini.

\subsection{Subbab Kedua}
\label{subsec:subbab-kedua}

% Tulis konten di sini.

%%%%%%%%%%%%%%%%%%%%%%%%%%%%%%%%%%%%%%%%%%%%%%%%%%%%%%%%%%%%%%%%%%%%%%%
% BAB 3 - METODOLOGI
%%%%%%%%%%%%%%%%%%%%%%%%%%%%%%%%%%%%%%%%%%%%%%%%%%%%%%%%%%%%%%%%%%%%%%%

\mychapter{3}{BAB 3 METODOLOGI}

% Gambar diagram alir metodologi (opsional)
% \begin{figure}[tph]
%   \centering
%   \includegraphics[width=.5\linewidth]{img/method.png}
%   \caption{Diagram Alir Metodologi}
%   \label{fig:diagram-alir}
% \end{figure}

\section{Metode Penelitian}
\label{sec:metode-penelitian}

% Tulis metode penelitian di sini.

\section{Tahapan Penelitian}
\label{sec:tahapan-penelitian}

% Tulis tahapan penelitian di sini.

% \include{bab4}
% \include{bab5}
% \include{bab6}
% \include{bab7}

% -----------------------------------------------------------------
% Bibliography
% -----------------------------------------------------------------
\printbibliography[title={DAFTAR PUSTAKA}]

\end{document}

%%% Local Variables:
%%% coding: utf-8
%%% mode: latex
%%% TeX-engine: xetex
%%% TeX-master: "skripsi"
%%% TeX-command-extra-options: "-shell-escape"
%%% ispell-local-dictionary: "id"
%%% End:
